\chapter{Document Authoring Guide}
\ChapterLabel{DocumentGuide}
\index{LaTeX}
\index{document conventions}
\index{build tools}
\index{AI agent!LaTeX installation prompt}
\index{Appendix!Document Authoring Guide}

\section{How to Use This Appendix}
Use this appendix to maintain authoring conventions, repository locations,
diagram sources, review rules, tool installation, and PDF build instructions.
Update it whenever the document structure or toolchain changes.

\section{Template Provenance}
\index{template provenance}
This template was developed by \textbf{Darian Muresan, Ph.D.}, in conjunction
with the AI agents \textbf{OpenAI Codex} and
\textbf{Anthropic Claude Code}. Questions, corrections, and proposed
improvements may be sent to
\href{mailto:dmuresan@stevens.edu}{dmuresan@stevens.edu},
\href{mailto:darian@dmmd.net}{darian@dmmd.net}, or
\href{mailto:darian@pictomic.com}{darian@pictomic.com}.

Keep the template attribution distinct from the project-specific document
owner and organization. A derived project may replace the project metadata
while retaining this provenance statement for the underlying template.

\section{AI-Assisted Authoring}
\index{AI agent!architecture workflow}
The authoritative instructions for AI-assisted architecture work, including
task routing, limited context loading, workflows, example prompts, completion
checks, and human-review boundaries, are contained in
\ChapterRef{AIAgentGuide}. This appendix retains installation, build, file, and
LaTeX authoring guidance that applies to both human authors and AI agents.

\section{File Naming}

Every chapter file uses the lowercase namespace \texttt{dsn}, followed by a
Pascal-style responsibility name and the \texttt{.tex} extension. Examples are
\texttt{dsnRequirements.tex}, \texttt{dsnTests.tex}, and
\texttt{dsnValidation.tex}. Supporting diagram files use the same prefix.

\section{One Responsibility per Chapter}

A chapter answers one major engineering question. When content belongs
to another responsibility, link to its authoritative chapter instead of
copying it. In particular, tests define procedures, verification records
requirement compliance, validation evaluates stakeholder use, and release
decides whether the evidence is sufficient.

\section{Adding a Chapter}

Create the namespaced file, begin it with a chapter title and stable chapter
label, add useful index terms, and include it once from
\texttt{dsnManual.tex}. Add or change artifact mappings only in the canonical
traceability matrix.

\begin{verbatim}
\chapter{Example Chapter}
\ChapterLabel{ExampleChapter}
\index{example}
\index{Chapter!Example Chapter}

Chapter content goes here.
\end{verbatim}

\section{Building the PDF}

The build files regenerate all Graphviz diagrams before compiling the document. They require:

\begin{itemize}
  \item Graphviz, including the \texttt{dot} command;
  \item a LaTeX distribution containing \texttt{pdflatex}, \texttt{bibtex}, and \texttt{latexmk}; and
  \item the \texttt{makeglossaries} and \texttt{makeindex} commands.
\end{itemize}

\subsection{Installing the Tools on macOS}

The full MacTeX distribution is recommended because it includes the LaTeX tools and packages used by this template. Complete these steps from an administrator account:

\begin{enumerate}
  \item Open Terminal and install Apple's command-line tools, which provide \texttt{make}:
\begin{verbatim}
xcode-select --install
\end{verbatim}
  Complete the installer dialog before continuing.

  \item Install Homebrew by following the current instructions at \url{https://brew.sh/}. Then install Graphviz and MacTeX:
\begin{verbatim}
brew install graphviz
brew install --cask mactex
\end{verbatim}
  MacTeX may instead be downloaded directly from \url{https://tug.org/mactex/}.

  \item Close and reopen Terminal so the MacTeX command-line tools are added to the system path.

  \item Verify that every required command is available:
\begin{verbatim}
command -v make dot pdflatex bibtex \
  makeglossaries makeindex latexmk
\end{verbatim}
  Each tool shall print a path. A missing line identifies the tool that still needs to be installed or added to the system path.
\end{enumerate}

\subsection{Installing the Tools on Windows}

Use a current Windows 10 or Windows 11 computer and complete these steps:

\begin{enumerate}
  \item Download and run the Basic MiKTeX Installer from \url{https://miktex.org/howto/install-miktex}. Choose a private per-user installation when appropriate and select \emph{Always install missing packages on-the-fly} for a computer with Internet access.

  \item Open MiKTeX Console, check for updates, and install them. If \texttt{latexmk} is not found during verification, use the Packages page in MiKTeX Console to install the \texttt{latexmk} package.

  \item Open PowerShell and install Graphviz with Windows Package Manager:
\begin{verbatim}
winget install graphviz
\end{verbatim}
  If \texttt{winget} is unavailable, use the current 64-bit installer from \url{https://graphviz.org/download/} and enable the option that adds Graphviz to the system path.

  \item Close and reopen PowerShell, then verify the required commands:
\begin{verbatim}
Get-Command dot,pdflatex,bibtex
Get-Command makeglossaries,makeindex,latexmk
\end{verbatim}
  Each tool shall produce a command record. A command-not-found error identifies a missing installation or system-path entry.
\end{enumerate}

On an organization-managed computer, authors shall ask the system administrator
before installing software or changing a managed execution policy.

\subsection{Installing the Tools on Ubuntu or Debian Linux}

Open a terminal and install Graphviz, GNU Make, Latexmk, and a TeX Live set
that contains the packages used by this template:

\begin{verbatim}
sudo apt update
sudo apt install graphviz make latexmk texlive-latex-extra \
  texlive-fonts-recommended texlive-bibtex-extra \
  texlive-extra-utils
\end{verbatim}

Verify the required commands:

\begin{verbatim}
command -v make dot pdflatex bibtex \
  makeglossaries makeindex latexmk
\end{verbatim}

For another Linux distribution, install the equivalent Graphviz, GNU Make,
Latexmk, and full or extended TeX Live packages with its package manager.

\subsection{AI-Agent Installation}
\label{Section:AIAgentInstallation}

An AI coding agent can inspect the host, identify missing commands, request
approval for system changes, install the appropriate packages, and verify the
build. Give the agent the project directory and the following reusable request:

\begin{verbatim}
In this LaTeX project, inspect the build files and determine
which operating system and package manager are available.
Verify make, dot, pdflatex, bibtex, makeglossaries,
makeindex, and latexmk. Install only missing tools using the
platform's standard package manager. Ask for approval before
administrator-level, system-wide, or networked installation.
Do not replace project sources or remove user files. After
installation, regenerate the Graphviz diagrams, run the full
platform build, inspect the log for errors, unresolved
references, and overfull boxes, and report the generated PDF.
\end{verbatim}

If the agent cannot install software, ask it to report the exact missing
command, the attempted verification, and platform-specific installation
instructions. Never give an agent an unrestricted instruction to delete build
directories or modify system security policy.

\subsection{Building on macOS and Linux}

From the template directory, run:

\begin{verbatim}
make -f makefileMacLinux pdf
\end{verbatim}

Use \texttt{make -f makefileMacLinux diagrams} to regenerate only the diagrams and \texttt{make -f makefileMacLinux clean} to remove generated build files while preserving authored sources and the final PDF.

\subsection{Building with Windows PowerShell}

From PowerShell in the template directory, run:

\begin{verbatim}
.\makefilePowerShell.ps1 -Action pdf
\end{verbatim}

Use \texttt{-Action diagrams} or \texttt{-Action clean} in place of \texttt{-Action pdf} to run those operations independently. If local execution policy prevents scripts from running, use an approved organizational policy or allow scripts only for the current PowerShell process:

\begin{verbatim}
Set-ExecutionPolicy -Scope Process Bypass
.\makefilePowerShell.ps1 -Action pdf
\end{verbatim}

\subsection{Manual Build Sequence}

When the platform build script is unavailable, run the following commands from
the template directory. The final two LaTeX passes resolve tables of contents,
citations, glossary links, the index, and cross-references.

\begin{verbatim}
dot -Gdpi=180 -Tpng diagrams/dsnVModel.dot \
  -o png/dsnVModel.png
dot -Gdpi=180 -Tpng diagrams/dsnUseCaseDiagram.dot \
  -o png/dsnUseCaseDiagram.png
pdflatex -interaction=nonstopmode -halt-on-error dsnManual.tex
bibtex dsnManual
makeglossaries dsnManual
makeindex dsnManual.idx
pdflatex -interaction=nonstopmode -halt-on-error dsnManual.tex
pdflatex -interaction=nonstopmode -halt-on-error dsnManual.tex
\end{verbatim}

The provided build scripts regenerate every \texttt{.dot} diagram, not only the
two shown in this abbreviated manual sequence.

\subsection{Build Verification and Troubleshooting}

Confirm that the output PDF opens, all expected chapters and appendices appear,
the glossary and index are populated, hyperlinks reach their targets, and no
table or figure is clipped. Search \texttt{dsnManual.log} for fatal errors,
undefined references, missing files, and overfull boxes. Common remedies are:

\begin{itemize}
\item open a new terminal after installation so its search path is refreshed;
\item on macOS, verify that \texttt{/Library/TeX/texbin} is on the search path;
\item rerun the complete build after adding citations, glossary entries, index
entries, figures, labels, or references;
\item regenerate diagrams after changing a \texttt{.dot} source; and
\item install the named LaTeX package when the log reports a missing
\texttt{.sty} file.
\end{itemize}

\section{Authoring Rules}

\begin{itemize}
\item Add the newest approved change to the top of the document history.
\item Introduce specialized terms through the glossary and use them consistently.
\item Give every controlled artifact a stable explicit ID.
\item Use \texttt{shall} only for binding requirements.
\item Keep normative requirements separate from explanation and examples.
\item Store editable diagram sources in \texttt{diagrams/} and rendered images in \texttt{png/}.
\item Compile and inspect the PDF after structural or table changes.
\end{itemize}

\section{Citations and Bibliography}

Store controlled reference records in \texttt{dsnReferences.bib}. Cite a
source from a chapter with a stable BibTeX key:

\begin{verbatim}
\cite{ISO29148}
\end{verbatim}

The main file applies the style and prints one document-level bibliography:

\begin{verbatim}
\bibliographystyle{IEEEtran}
\bibliography{dsnReferences}
\end{verbatim}

Both platform build scripts run BibTeX automatically.

Bibliography entries support claims and definitions; they do not replace
project-owned requirements, architecture decisions, test procedures, or
approval evidence.

\section{Selective Builds}

For a large project, LaTeX's \verb|\includeonly| command may select chapters
during drafting. The main file contains a commented example:

\begin{verbatim}
% \includeonly{
%   prologue,
%   dsnRequirements,
%   dsnTests
% }
\end{verbatim}

Remove the leading comment characters and list include names without the
\texttt{.tex} extension. LaTeX can retain cross-references and citations from
omitted chapters by reading their existing auxiliary files, but those files
may be absent or stale. Selective builds are therefore drafting aids only.
Hyperref may also report missing destinations for links into deliberately
omitted chapters; these draft warnings shall disappear in the full build.
Comment out the complete \verb|\includeonly| block and run a full build before
verification or release so missing chapters, references, glossary entries,
citations, and evidence cannot be hidden.
